We start our story with one of the primordial sciences: Astronomy. Two infamous personalities occupy the setting for the first concepts we discuss, namely the astronomers Tycho Brahe $(1546-1601)$ from Denmark and Johannes Kepler $(1571-1630)$ from Germany.  Tycho was perhaps the best observational astronomer we ever knew of, his observations were precise and all compassing. His obsession and careful noting with the motion of the sky objects is unheard. Kepler has a more mathematical approach to astronomy, there is a reason that to think of Kepler is usually of think of his mathematical laws for orbits. Tycho actually met Kepler and asked him to work as his assistant, unfortunately, not much time passed between this event and Tycho's death. However, the incredible amount of observations made during his lifetime would go on to empower Kepler to correctly deduct and improve his theories. Newton $(1643-1727)$, who doesn't require introductions as one the most prolific and successful physicist, astronomer and mathematician, came later to derive by first principles and the laws of gravitation the observational theory. By this point, the motion of planets was well understood and there was a good idea on how to work with other heavenly bodies.

There was some difficulty, however, when it came to the motion of comets and like bodies from outer space. The observations were scarce and imprecise, these bodies made appearance once to only a few times if one were lucky enough find it. The theory of gravitational laws appointed ellipsis and parabolas as contenders, but the observational data were still lacking. How to, from a small set of error-prone, derive the original curves of motion; the problem is better stated by Gauss $(1777-1855)$. 

\begin{citacao}
	To determine the orbit of a heavenly body, without any hypothetical assumption,
	from observations not embracing a great period of time, and not allowing a selection
	with a view to the application of special methods.  \cite{gauss1809theoria}
\end{citacao}

It was Gauss who solved the problem. From our current understanding, Gauss used a type of \textit{Maximum Likelihood Estimation} method, that is, he tried to derive parameters of the underlying curve from a set of observational points without any previous knowledge or assumption. The idea is to try to minimize some error function that is correlated with the distance of the observational points to the predicted curve. One might remember for example of the \textit{least squares} method. As mentioned, observational data carried errors, to perform such methods Gauss would need to understand how to weight the distance of the observational point to the theoretical curve. Surely, points very far from a predicted curve should be few, this is assumption one: The error function must be decaying. Another assumption was that this error function was symmetric, an even function. These assumptions led Gauss to the normal curve, or what it sometimes called the Gaussian distribution. 

\begin{figure}[h!]
	\centering
	\begin{tikzpicture}
		\begin{axis}[
			xmin=-3, xmax=3,
			width=\textwidth, height=6.0cm,
			%xticklabel=\empty,
			xticklabels={-3, , -2, , -1, , 0, , 1, , 2, , 3},
			ymin=0, ymax=0.55
			]
			\addplot[name path=f1,solid, black, domain = -3:3, samples = 100] {(e^(-x^2/2))/sqrt(2*pi)};
			\path[name path=axis] (axis cs:-3,0) -- (axis cs:3,0);
			\tikzfillbetween[of=f1 and axis]{pattern=dots};
		\end{axis}
	\end{tikzpicture}
	\caption{Plot of $e^{(-x^2)/2} / \sqrt{2\pi}$, the Gaussian distribution.}
\end{figure}

Today, it is well-known among physicists and mathematicians alike that random independent phenomena share common characteristics when dealt in large quantities. The most important and known example is surely the convergence to the normal curve. Supposing you're working under some reasonable assumptions on the random variable at hand, if one looks at the distribution of the mean of such random variables when the number of samples increases one should see convergence to a known curve: The normal. The most important assumptions here is that you're dealing with a distribution with a somehow contained tail - finite moments, and independent samples. This may be the core result in much of applied statistical research. This is formally expressed in terms of the central limit theorem: Take $\{x_n\}_{n\geq1}$ independent and identically distributed random variables, if we center and re-scale the variables $x_n \rightarrow y_n \deff (x_n - \E(x_n)) / \sqrt{\Var(x_n)}$, then we have the limit
\begin{equation}
	\lim_{n\rightarrow \infty} \p\left( \frac{\sum_{k=1}^{n} y_k}{\sqrt{n}} \leq t \right) =F^{(G)}(t; 0, 1)
\end{equation}
where $F^{(G)}(t; 0, 1)$ is simply the cumulative distribution for the normal distribution. This in an amazing result as it appoints the Gaussian distribution as a universal distribution for a large number of uncorrelated random processes. 

There is however, many phenomena where these assumptions can't hold. One can think for example of an interacting system of particles as a highly important example of correlated phenomena. For these cases another universality result must be considered, one where the Gaussian distribution is substituted with a more appropriate distribution. An important historical example appeared in physics in the last century. When working with heavy nuclei, physicist around the $1950$'s struggled a lot to properly understand what was going on with the system compared with the absolute success of quantum mechanics in simpler nuclei such as hydrogen. This type of work was condemned to be impossible: Theoretically, these atoms are huge and the interactions and intern structures were barely known. Dyson $(1923 - 2020)$, a British-American mathematician and physicist,  described the situation as follows. 

\begin{citacao}
	We picture a complex nucleus as a ‘black box’ in which a large number of particles are interacting according to unknown laws. The problem then is to define in a mathematically precise way an ensemble of systems in which all possible laws of interaction are equally probable. \cite{Dyson1962es}
\end{citacao}

Physically, to understand a system one might want to write out its Hamiltonian - in a way, a expression of its energy. This takes the form of a matrix with size of the order of the independent agents you have in the system. For large nuclei, you would deal with huge matrices.  Computationally, those were also useless, computers were not, and were not expected to be, at a point of evaluating meaningfully, even if one understood well the system, the evolution of such system. That was a typical experiment that was done to study properties of such nuclei and that physicists were very interested in understanding theoretically the results. The nuclei scatter is a technique of projecting neutrons in nuclei at multiple energies to detect resonance energies. These energies are related to the internal energy and of much use to researches. In a conference on nuclear physics in $1956$, Eugene Wigner $(1902-1995)$, a Hungarian-American theoretical physicist,  presented his work regarding the theoretical statistics of the spacing distribution of adjacent neutron resonances in heavy nuclei. He said the following.

\begin{citacao}[english]
	Perhaps I am now too courageous when I try to guess the distribution of the distances between successive levels (of energies of heavy nuclei). Theoretically, the situation is quite simple if one attacks the problem in a simpleminded fashion. The question is simply what are the distances of the characteristic values of a symmetric matrix with random coefficients. (WIGNER, 1956)
\end{citacao}

The fact is that Wigner had had a breakthrough that would change of correlated systems and introduce \sigla{RMT}{Random Matrix Theory}. He found a way to model scattering energies as eigenvalues of a specific model of random matrix - though of simply as a matrix of random entries with some given distribution, correlated or not. To understand how and what, we must think of what it means to model a system by a random matrix.  For that, Deift $(1945)$, a South-African mathematician, has a very nice analogy.

\begin{citacao}[english]
	This procedure can  be viewed as follows: A scientist wishes to investigate some statistical phenomenon. What he has at hand is a microscope and a handbook of matrix ensembles. The data $\{a_k\}$ are embedded on a slide which can be inserted into the microscope. The only freedom that the scientist has is to center the slide, $a_k \rightarrow a_k - A$, and then adjust the focus $a_k - A \rightarrow \tilde{a}_k = \gamma_a(a_k - A)$ so that on average one data point $\tilde{a}_k$ appears per unit length on the slide. At that point the scientist takes out his handbook, and then tries to match the statistics of the $\tilde{a}_k$’s with those of the eigenvalues of some ensemble. If the fit is good, the scientist then says that the system is well-modeled by RMT. \cite{deift2006universalitymathematicalphysicalsystems}
\end{citacao}

Wigner understood that the energies he was studying had the same statistic as the eigenvalues of random matrices from the so called \textit{Gaussian Ensemble} by the procedure described by Deift. This ensemble is nothing more than the set of hermitian variables with independent random centralized and normalized Gaussian variables. The distribution of eigenvalues is given by the famous Wigner Semi-Circle.
\begin{figure}[h!]
	\centering
	\begin{tikzpicture}
		\begin{axis}[
			xmin=-1.5, xmax=1.5,
			width=\textwidth, height=7.0cm,
			ymin=0, ymax=0.8,
			xtick distance=0.5
			]
			\addplot[name path=f1, solid, black, thin] table [col sep=comma] {./plotdata/SemiCircle.csv};
			\path[name path=axis] (axis cs:-3,0) -- (axis cs:3,0);
			\tikzfillbetween[of=f1 and axis]{pattern=dots};
		\end{axis}
	\end{tikzpicture}
	\caption{Plot of the Semi-Circle.}
\end{figure}

 In the bulk, the distribution can be described by the so-called Sine Kernel, the same statistic found in the atomic energies described by Wigner. This distribution also have other interesting kinks. For some point $e$ in the edge of the distribution
\begin{equation}
	\p\left(e + \frac{x}{c N^{2/3}} \leq t \right) = F(x)
\end{equation}
we have a special joint probability density function: the Tracy-Widom. 

A very classical model for these distributions is the problem of the maximal increasing sub-sequence, know by the name of one of its researcher, the Polish mathematician Stanislaw Ulam $(1909 - 1984)$, as \textit{Ulam's problem} \cite{Ulam}. Suppose you take a sequence of integer number and shuffle them in what we call a permutation. There are many possible permutations, and we will consider for this application that they are all equally likely. Naturally one can get for example
\begin{equation*}
	\pi = 54172386
\end{equation*}
as a permutation of the numbers $\{1, 2, 3, 4, 5, 6, 7, 8\}$. One can notice that the number are not necessarily increasing or decreasing globally, but that it may be the case locally and even more the case if one is willing to ignore the non-conforming numbers. That is, take for example the sub-sequence of $\pi$ given by $1236$. Note that these numbers are not sequential but nevertheless are an increasing sub-sequence of the original permutation. Suppose that we consider the distribution of sub-sequences lengths. For many reasons one might be interested in the maximal increasing sub-sequence $l_N$. Of course, calculating $l_N(\pi)$, for $\pi$, is a trivial exercise and for larger sequences, only requires patience and will force. However, in most application one wants to understand the typical size of such maximal large sequences, i.e., what size one might expect to get from a random sampled permutation of the numbers. Some direct assumptions can be made: this size must depend on a parameter $n$, the size of the initial sequence. Of course, it ca not be bigger than $n$ or smaller than $1$; however it stands to reason that this size must have some scaling with $n$.

This problem is also highly linked with the problem of the \textit{matching sequences} \cite{majumdar2007randommatricesulamproblem}, or more descriptively, the problem of the longest common sub-sequence. We deal again with a permutation of some alphabet of symbols, before we used the integers only allowing each number to appear once. Now, we could still use numbers but for a non-suspicious motive we will use the four letter alphabet $ACTG$ and allow repetition of symbols in a sequence. Consider for example the sequences $\alpha, \beta$ given by
\begin{align*}
	\alpha \quad & : \qquad A \quad C \quad G \quad C \quad T \quad A \quad C \\
	\beta \quad & : \qquad C \quad T \quad G \quad A \quad C
\end{align*}
Again, we want to look for, non necessarily adjacent, sub-sequences of symbols. For example, a sub-sequence of $\alpha$ could be $CGAC$, corresponding to the second, third, sixth and seventh letters. As it so happens, we chose a sub-sequence that is common to both $\alpha$ and $\beta$. Moreover, as $\beta$ is of size $5$, this might just be the longest (one of) matching sequences as long as $\beta$ is not contained as a sub-sequence in $\alpha$. For the biology interested reader, an application of this type of metric may have become clear. The fact is that DNA are sequences of a four lettered alphabet such as $\alpha$ and $\beta$. It is our understanding that genes responsible for specific proteins evolve with time and that, finding this longest common sequences between genes of different species may indicate what is conserved though evolution. One may also be interested in checking that two genes are the same, or if they are considerably different. Just by chance, of course, one must expect the longest matching sub-sequence to be relatively large for a few-lettered alphabet. However, if one understands whats the expectation of such length and how lengths deviates from this mean, it gets considerably easier to make statements that two sequences are evolutionary close or far apart.

Note that in both cases these variables can not be understood as independent random variables and that the classical central limit theorem is useless for its study in the enunciated form. For the longest increasing subsequence problem, historically, \citeonline{Ulam}, discussed the idea of using Monte Carlo approaches to study the asymptotic behavior of $\E[l_N]$, the expectation of the longest increasing sequence length when $N \rightarrow \infty$. Based on his simulations, Ulam conjectured that it exists a constant $c$ such that $\E[l_N] = \sqrt{N}c$ when $N$ went to infinity. Later on, \citeonline{Hammersley1972AFS}, introduced a Poissonized version of this problem and formally proved that the constant $c$ exists. Five years afterwards, \citeonline{LOGAN1977206} and, independently, \citeonline{VerKer77} showed that $c = 2$. Much later, \citeonline{odlyzko_2000_yd8at-ym210}, performed analytical simulations indicating that 
\begin{equation}
	\lim_{N \rightarrow \infty} \frac{\Var[l_N]}{N^{\frac{1}{3}}} = c_0 \equiv 0,819\cdots
\end{equation}
and also that
\begin{equation}
	\lim_{N \rightarrow \infty} \frac{\E[l_N] - 2\sqrt{N}}{N^{\frac{1}{6}}} = c_1 \equiv -1.748\cdots.
\end{equation}
At this point, the expectation of $l_N$ was well known. However, the proper variable had not been taken face on. Finally, Ulam's problem was solved by \citeonline{baik1999distributionlengthlongestincreasing}, who studied the convergence in distribution of the re-scaled variable $ \chi_N \deff (l_N - 2\sqrt{N})/ N^{\frac{1}{6}}$ and discovered that it converges in distribution to $\chi$: 
\begin{equation}
	\lim_{N \rightarrow \infty} \p\left( \frac{l_N - 2\sqrt{N}}{N^{\frac{1}{6}}} \leq x\right) = F_{(2)}(x).
\end{equation} 
\begin{figure}[h!]
	\centering
	\begin{tikzpicture}
		\begin{axis}[
			xmin=-5, xmax=3,
			width=\textwidth, height=6.0cm,
			xticklabels={ , , , -4, , , , -2, , , , 0, , , , 2, ,},
			ymin=0, ymax=0.55,
			]
			\addplot[name path=f1, solid, black, thin] table [col sep=comma] {./plotdata/TracyWidom-1.csv};
			\path[name path=axis] (axis cs:-5,0) -- (axis cs:3,0);
			\tikzfillbetween[of=f1 and axis]{pattern=dots};
		\end{axis}
	\end{tikzpicture}
	\caption{Plot of $F_{(2)}(x)$.}
\end{figure}

The daunting resemblance to the central limit theorem should be noted. This cumulative distribution function $F_{(2)}(x)$ ended up being know now as the \textit{Tracy-Widom}. By good measure we must also state a result from the matching sequence example. The result is only valid for a simplified model: The \textit{Bernoulli matching model}; got if one ignores some correlation that appears in the modeling. This is not completely bad as these correlations tend to zero when the number of letters in the alphabet increases. Name the size of the longest matching sequence $L^{(BM)}_{i,j}$ where $i,j$ are the size of the first and second sequences used. If $i=j=n$, for large $n$, one gets that
\begin{equation}
	L^{(BM)}_{n,n} \rightarrow \gamma^{(BM)}_c n + f(c) n^{1/3} \chi	
\end{equation}
where $\gamma_c^{(BM)} = 2/(1+\sqrt{c})$, $\chi$ is a random variable with a $n$-independent distribution and $\p(\chi \leq x) = F_2(x)$ is precisely the Tracy-Widom distribution, again.

This special function appears many other times in many other different models and phenomena. It appears to be related to the universality of phase transitions - such as the physical transition of water or superconductivity. One of the indications of this development was the work of Majumdar and Schehr that shown that for large $n$, the Tracy-Widom to describe the crossover function between stable and unstable regimes of a random dynamical system - a phase transition, such as the ones observed in Coulomb gases or two-dimensional QCD \cite{Kieburg_2014}. Eventually physicists realized that the energy curves of certain strongly correlated systems have a kink at $\sqrt{2N}$. The associated peak for these systems is the Tracy-Widom distribution, which appears in the third derivative of the energy curve - that is, the rate of change of the rate of change of the energy’s rate of change. 

More than that, for a large class of random matrices called Wigner Matrices, independent of individual probability density of matrices, the eigenvalue distribution tends to the Wigner's Semi-Circle. It represents a class of universality in matrices. Interestingly enough, the eigenvalue dynamics can be often interpret as a physical gas of charged interacting particles submitted to some confining external potential know as a Coulomb gas. Note that such particles would repel each other, however unnatural this might sound for eigenvalues of random matrices, this is what is observed: Eigenvalues configurations with two extremely close eigenvalues are much less likely. This creates a bridge to a theory of dynamical systems and point processes that lead us to consider what happens when such gases are taken to the thermodynamic limit. That is, what happens when we increase the number of charged particles, the mass number of an atomic nuclei or even the length of the sequences considered before. This give us under some appropriate considerations, to asymptotic equilibrium configurations of the points and is a major field of study in RMT. It also motivates ways of simulating this equilibrium measures, namely, using usual particle dynamics methods such as in \citeonline{Chafal}.

\begin{figure}[h!]
	\centering
	\begin{tikzpicture}
		\begin{axis}[axis line style={opacity=0}, ticks=none,
			xmin=-1.5, xmax=1.5,
			width=8cm, height=7.8cm,
			ymin=-1.5, ymax=1.5,
			xticklabel=\empty,
			yticklabel=\empty,
			scatter/classes={%
				a={ mark size=0.5pt}}
			]
			\addplot[skip coords between index={100}{500}, opacity=1, scatter, mark=x, black, only marks] table [col sep=comma] {./plotdata/Symmetric/5lines/a3t1.csv};
			\addplot[opacity=0.05, scatter, mark=*,draw opacity=0,black, only marks] table [col sep=comma] {./plotdata/Symmetric/a3t1.csv};
		\end{axis}
	\end{tikzpicture}
	\caption{Plot of a simulation of a Coulomb gas. In black crosses we plotted one sample from the simulation. In shadowy dots, any \enquote{averaged} samples indicating the equilibrium region.}
\end{figure}

In recapitulation, Wigner noted that modulo some symmetry considerations, all nuclei in the same symmetry class exhibited universal behavior; and that in all cases, the energies exhibited repulsion, or, more precisely, the probability that two energy levels would be close together was small.  If one equips the space of symmetric matrices with a probability measure of smooth density, the probability of a matrix $\m{M}$ having equal eigenvalues would be zero and the eigenvalues $\lambda_1, \cdots , \lambda_n$ of $\m{M}$ would be a random set with repulsion already built in. With this, Wigner had a model with repulsion. But why choose the Gaussian Orthogonal Ensemble, composed of real Gaussian random matrices? This is where the universality constraint comes into play. We quote from Wigner, where he mentioned the previously used matrices defined by \citeonline{Wishart1928-hi} to study covariance matrices in real world data.
\begin{citacao}
	“Let me say only one more word. It is very likely that the curve in [Figure 1] (referring to the Semi-Circle) is a universal function. In other words, it doesn’t depend on the details of the model with which you are working. There is one particular model in which the probability of the energy levels can be written down exactly. I mentioned this distribution already in Gatlinburg. It is called the Wishart distribution. ...” 
\end{citacao}

So in this way Wigner introduced a model with built-in repulsion and in which the energy level distribution could be calculated explicitly. 

This type of approach became better know and phenomena in the same universality class started to pup-up many times. A well know story regard H. Montgomery (1944), an American mathematician, and Dyson. Montgomery works with analytical number theory and is better know for his \textit{pair correlation conjecture} on the zeros of the Riemann zeta function $\zeta(s)$ - perhaps one of the most know analysis open problem today. The idea is that Montgomery found a way to re-scale the non-trivial zeros of the Riemann zeta function and find a pair correlation between these objects. That is, he found some statistical information on the distribution of such zeros, a very hard problem. During a visit in the $1970$'s to the Institute for Advanced Study in Princeton he was suggested to present his result to Dyson. Before Montgomery could state his result however, Dyson took out a pen and predicted some function that looked like
\begin{equation}
	R(a, b) = \int_{a}^{b} 1 - \left( \frac{\sin(\pi r)}{\pi r}\right)^2 \dd r.
\end{equation}
Of course, Montgomery should have felt nothing short of deep confusion as how his result should be that obvious for Dyson. The explanation for such a guess is that Dyson believed that the zeros of the zeta function might behave like the eigenvalues of a random Gaussian Unitary Ensemble (GUE) matrix (composed of unitary matrix with Gaussian entries), and, were that the case, this would be exactly the formula for the two-point correlation function. This was derived from the well know by the time, Sine Correlation Kernel.

This of course, was incredibly interesting and was looked for in all sorts of phenomena. Another interesting result is the work of \citeonline{Abul_Magd_2006}, that investigated the gap distances between parked cars to find that these distances behaved as the distances of eigenvalues for the GUE. The explanation given was that drives tended to park near each other but steel tried to keep a reasonable distance to maneuver safely, in other words, they behaved just like a charged particle gas with a confining potential! Another instance of this type of result was in the $1990$'s when two Czech Physicists took interested in the transportation system of the city of Cuernavaca, in Mexico. In this work, \citeonline{Krb_lek_2000} suspected that the peculiarities of the system public buses could result in a result expressed in terms of know objects in RMT. There were no public buses in the city, buses were private and did their own timetables only following the approved routes designed by the city. As there were no coordination for the buses times, sometimes buses were right before another line and could not take anyone. To avoid that, the city installed sensors and started to sell the information about buses leaving the stations. Now, with theses information, buses started to avoid each other but of course, still wanted to to the route as soon as possible to avoid loosing clients. This resulted in the same behavior noted before: buses were behaving as the aforementioned gas!

This modeling was repeated for many other highly correlated processes and was very successful as a general theory of random correlated systems. This can now be observed in the plenitude of posterior system modeled with such statistics, for example: the aforementioned Ulam's Problem \cite{Ulam} and Matching Random Sub-sequences \cite{majumdar2007randommatricesulamproblem} but also in the Zeros of the Riemann Zeta Function \cite{Montgomery}, Tiling problems (Aztec Diamond for instance) and Interacting Random Walks, to name a few. 

\section{Overview of the dissertation}

Our main goal with this work is to obtain, following the recent results of \citeonline{Byun_2023}, asymptotic expressions, on the number of particles, for the normalizing function - called partition function - of the distribution of eigenvalues for the ensemble of matrices we denominate planar ensembles. There is much to unpack in this sentence. Namely, one can mention that such distributions of eigenvalues behaves as a fluid gas of interacting particles under some external potential; we call such system a Coulomb gas. Such gases carry some physical implications that we make good use in the mathematical theory. One can imagine that if the potential is confining, with the increase of particles we would get the thermodynamic limit of the system, i.e., particles should stabilize on some equilibrium distribution that minimizes energy, or equivalently maximizes entropy. Those will be our equilibrium measures for the eigenvalues. Furthermore, these systems behave as a point dynamic called \sigla{DDP}{determinantal point process}, i.e., the joint distribution of the points can be expressed as a determinant of a function called correlation kernel. These kernels are expressed as a functions of special monic polynomials that are orthogonal in respect to a so-called weight function that depends on the confining potential on the gas system. Now, the partition function is simply the summation overall possible micro-states of the system. This can be expressed as an integral of the eigenvalues density over the possible configurations and therefore as a function of the same correlation kernel as before. To perform the asymptotic for the partition function is then to perform the asymptotic for the orthogonal polynomials or its norming constants in a sense. Luckily the integrals we need to perform asymptotic on are of the physics favorite form, the integrand is the exponential of a function with one or a few isolated maxima points multiplied by some increasing parameter. This is good because these integrals can be, given some conditions, approximated asymptotically by a Gaussian integral and some error bound. This asymptotic results are extremely interesting as the coefficients on each order of $n$, the size of the system, have physical interpretations such an internal energy, entropy and others less straightforward. 

On \autoref{chapter: RMT} we will introduce our main culprits: The Random Matrices in \sigla{RMT}{Random Matrix Theory}. Our goal with this chapter is to formalize notation and language for the rest of the work while introducing the main ideas one need to understand when dealing with RMT. Random matrices can be of many types, we will divide the types of ensembles and introduce some that are closer to our work or that are central to the theory. Dealing with random matrices one must also decide on with quantities are of interest in respect to this objects, it may be the expected characteristic polynomial, the trace, the distribution of eigenvector or eigenvalues, among others. The ideal of orthogonal polynomials and how they fit into the story of random matrices are also briefly touched as needed for the development. This chapter in essential in understanding what goes on in \autoref{chapter: equilibrium} and \autoref{chapter: planarensembles} and important to be connect with the ideas developed in \autoref{chapter: pointprocesses}.

On \autoref{chapter: pointprocesses} we explored the mathematical structure of the point dynamics we gain when dealing with eigenvalue statistics: Point Process (PP) and more specifically for us, Determinantal Point Process (DDP). On this chapter we reinterpret some quantities taken from the random matrices as function of the dynamic itself. This allow us to better understand what will be called the $k$-moments and correlation functions. This also lead us to the main structure in the distribution of eigenvalues of invariant ensembles; the determinant of the correlation kernel. We also show some properties for such kernels and explore a very important ideal of transformations and limit of these object. This connects nicely with the ideal of universality in RMT and the different local regime found in eigenvalue distributions, here represented by different scaling and transformations. This chapter is a good complement of \autoref{chapter: RMT} and provides some theoretical motivation for what comes in \autoref{chapter: equilibrium} and \autoref{chapter: planarensembles}.  

On \autoref{chapter:asymptotics} we construct the necessary and appropriate background in asymptotic methods for integral functions, to be used when performing asymptotics of orthogonal polynomials and related objects. There are two separate ideas in this chapter. The first, about quadrature methods, is an excursion and done as technical requirement for \autoref{chapter: planarensembles}. This works on a way to approximate, asymptotically, integrals as summations using the Euler-Mclaurin formula. The second part refers to the main ideas of asymptotics we use: the Laplace technique. This is a way to approximates integrals in which a major part of the contribution comes from a few isolated points in the domain.  This chapter is a bit out of place but a strict requirement for \autoref{chapter: planarensembles} and a soft one for a few ideas in \autoref{chapter: equilibrium}.

On \autoref{chapter: equilibrium} we more directly discuss the idea of Coulomb gases and how this idea lead us to equilibrium distributions for the eigenvalues. We begin by specifying the measure of Coulomb gases and noting that it is the same of invariant ensembles of matrices. Then we discuss quantities such as the energy of such system with and without an external confining field and how each case lead us to a minimization problem for the energy. The main results here refers to the Variational Equations that determine some techniques for identifying the equilibrium measure. This is also where we get to an generic expression for the partition function and where we motivate the asymptotic we will get on \autoref{chapter: planarensembles}.  This chapter is better read after \autoref{chapter: RMT}. It server as motivation and justification for equilibrium measures and the search for the expression on \autoref{chapter: planarensembles}.

On \autoref{chapter: planarensembles} we finally get to our main result: We prove the asymptotic expression of the partition function for some random normal ensembles. This is done following the work of \citeonline{Byun_2023}. Firstly we state our model and discuss the characterization of the results we want to get. The following pages are dedicated to a dense computation of the results for the two relevant cases of equilibrium measures: with disk and wing shaped supports. This is the culmination of the theoretical work done on the previously chapters and the main point of the work.

I hope you have a good journey into this work and enjoy the details and cool examples as much as I did. Good read!




